\section{Most failures are a false ``done''}
\label{sec:diag}

\paragraph{Taxonomy.} Each episode gets one primary category from fixed rules over its actions, screenshots and final score. \emph{False success}: the agent calls \texttt{answer} on a feasible task scored 0, or claims completion on an infeasible task. \emph{Loop}: at least three consecutive identical actions while less than 1\% of pixels change. \emph{Step limit}: the 50-step budget is exhausted without a loop. At step level, \fsp{} marks the final \texttt{answer} of a false success and \fsm{} that of a true one. A correctly worded refusal of an infeasible task is excluded from both. The rule file is hashed before any direction is fitted (current hash \texttt{\oswRulesSha}); later changes are logged with their hash and none altered a label. % source: JOURNAL.md "Labels figés avant toute AUROC OSWorld" and "Règle de label ajoutée avant toute AUROC"
Before fitting any direction we also require that the target category cover at least 25\% of failures and provide at least 40 positive steps (the GO rule).

\paragraph{Result.} On the diagnostic split (\oswNTasks{} tasks, \oswNScored{} scored), Holo-3.1-4B scores \oswMeanScorePct\% (the official OSWorld score, partial credit included, the convention used for every score in this paper) and fully solves \oswNSuccess{} tasks: \oswSuccessPct\%, 95\% Wilson CI \oswSuccessLo--\oswSuccessHi. An independent reproduction with the same public harness reports about 45.5\% at 100 steps on a different task set \citep{realtimegui}; we do not claim the two numbers are comparable. % source: JOURNAL.md "Baseline OSWorld 4B" (RealtimeGUI ~45.5 %)
Of the \oswNFail{} failures (Figure~\ref{fig:osw}), \oswCatFalseSuccess{} (\oswFsPct\%) are false successes, \oswCatStepLimit{} (\oswStepPct\%) have the step limit as primary category (\oswStepAny, or \oswStepAnyPct\%, exhaust the 50 steps at all, loops included), \oswCatLoop{} (\oswLoopPct\%) loop, and \oswCatOther{} is unclassified. \oswFsInfeasible{} false successes are claims of completion on infeasible tasks, and \oswFsFeasible{} are feasible tasks done wrongly or partially. On an infeasible task the correct action is \texttt{answer} with a refusal, so $P(\texttt{answer})$ is the wrong metric there; one such context, Chrome \texttt{480bcfea} at step 4, enters several screens below, and Appendix~\ref{app:infeasible} gives every count with and without it. These are mostly not misworded refusals: among the \infAllowed{} tasks open to construction, \infTasks{} are infeasible (\infEpisodes{} episodes), and of the \infAnswered{} final answers \infClaim{} claim success, while only \infInabilityTasks{} task states its inability, in words the harness rule misses. % source: data/pod_astra/infeas/reviewed_inventory.json, report.md
They are not early quits: the median false success ends after \oswFsMedianSteps{} steps. Every domain has false successes except Thunderbird, and Chrome, Calc, GIMP and VLC have four or more each (Fig.~\ref{fig:osw}a). % source: generated/tab_osw_domains.tex
False success passes the share criterion of the GO rule, but its \oswLabFsPos{} \fsp{} steps alone do not pass the count criterion. With a second pass over the same tasks, stopped at \rtmPassTwoScored{} scored tasks when the GPU returned to production, \bAdmFsPos{} \fsp{} and \bAdmFsNeg{} \fsm{} steps are admissible, above the required 40 (as of \bStatusUtc{} UTC).

\placefig{fig-osw}

\subsection{The model says ``done'' on false and true completions alike}
\label{sec:think}
We re-sampled the original model on six screening contexts from five tasks, two per stratum ($k=4$, production sampling, thinking enabled, invalid outputs counted; Fig.~\ref{fig:teaser}c and Table~\ref{tab:causal}, first row). It calls \texttt{answer} on \bPilotFsAnswer/\bPilotFsN{} false-completion samples and \bPilotTrueAnswer/\bPilotTrueN{} true ones, and never on mid-episode contexts (\bPilotHealthyAnswer/\bPilotHealthyN). It reproduces the terminal versus non-terminal distinction exactly; the false versus true distinction is invisible. Any repair must therefore change a decision the model takes with near certainty in both cases, and only in one of them.

The reasoning has the same shape in both cases. On \thN{} re-sampled reasoning traces at completion steps of \thNTasks{} construction tasks, \thFsPosN{} on false and \thFsNegN{} on true completions, every trace opens with a success claim (``The task has been completed successfully'', ``The document has been saved successfully'') at generated-token index \thTokMin--\thTokMax, counted from zero, with no verification word before it (\thNVerif/\thN), then lists the steps taken (median \thMedLen{} tokens). % source: data/pod_astra/think/direction_summary.json (verification_window empty for all events)
Because true completions show it too, this is a property of Holo's output format at completion steps, not a marker of false success. It locates the written conclusion, not the moment the decision forms inside the model. A direction fitted on these tokens separates the two pairs perfectly, but so do context length and a single lexical cue, so we give it no weight; steering along it is negative (Table~\ref{tab:causal}). What the format skips is the check: the Spotify episode states that the installation cannot finish and claims success one step later (Fig.~\ref{fig:teaser}a).
